\documentclass[12pt,a4paper,UTF8]{ctexart}
\usepackage[a4paper,top=22mm,bottom=22mm,left=22mm,right=22mm,headheight=15pt]{geometry}
\usepackage{amsmath,amssymb,mathtools,bm}
\usepackage{graphicx,xcolor,booktabs,longtable,array,enumitem,fancyhdr,caption,float,microtype,setspace}
\newcommand{\Needspace}[1]{\par\begingroup\dimen0=#1\relax\dimen2=\pagegoal\advance\dimen2 by-\pagetotal\ifdim\dimen0>\dimen2\newpage\fi\endgroup}
\newcommand{\needspace}[1]{\par\begingroup\skip0=\lastskip\vskip-\skip0\dimen0=#1\relax\advance\dimen0 by\skip0\vskip\dimen0\penalty-100\vskip-\dimen0\vskip\skip0\penalty9999\endgroup}
\usepackage[hidelinks,unicode]{hyperref}
\usepackage[most]{tcolorbox}
\definecolor{ContestBlue}{HTML}{17365D}
\hypersetup{pdftitle={第 048 套},pdfauthor={}}
\setstretch{1.24}
\setlength{\parindent}{2em}
\setlength{\parskip}{0.35em}
\setlength{\emergencystretch}{2em}
\setlist{itemsep=0.35em,topsep=0.4em}
\captionsetup{font=small,labelfont=bf,justification=centering}
\setcounter{secnumdepth}{0}
\pagestyle{fancy}
\fancyhf{}
\fancyhead[L]{\small 第 048 套}
\fancyfoot[C]{\thepage}
\renewcommand{\headrulewidth}{0.4pt}
\ctexset{
  section={format=\Large\bfseries\color{ContestBlue},beforeskip=1.2em,afterskip=0.7em},
  subsection={format=\large\bfseries\color{ContestBlue},beforeskip=1.0em,afterskip=0.5em},
  subsubsection={format=\normalsize\bfseries,beforeskip=0.8em,afterskip=0.35em}
}
\makeatletter
\newcommand{\tightlist}{\setlength{\itemsep}{0pt}\setlength{\parskip}{0pt}}
\newcommand{\maxwidth}{\ifdim\Gin@nat@width>\linewidth\linewidth\else\Gin@nat@width\fi}
\newcommand{\maxheight}{\ifdim\Gin@nat@height>0.72\textheight 0.72\textheight\else\Gin@nat@height\fi}
\newcommand{\pandocbounded}[1]{\begingroup\setkeys{Gin}{width=\maxwidth,height=\maxheight,keepaspectratio}#1\endgroup}
\makeatother
\begin{document}
\begin{titlepage}
  \centering
  \vspace*{58mm}
  {\Huge\bfseries\color{ContestBlue} 第 048 套\par}
  \vfill
  \vspace*{28mm}
\end{titlepage}
\section{第 1 题}\label{ux7b2c-1-ux9898}

车轮绕中心点 \(A\) 转动，连杆 \(OA\) 与主销轴 \(OB\)
始终垂直，整个轮架又可绕 \(OB\) 转动。

（1）从汽车正前方看，\(OB\) 与竖直方向的夹角为主销内倾角 \(\beta\)。已知
\(OA=l\)，四个车轮受到的支持力相同，汽车重力为
\(G\)。当轮架相对正中位置绕 \(OB\) 转过 \(\theta\)
时，求重力造成的自动回正力矩。

（2）从汽车侧面看，\(OB\) 与竖直方向有很小的主销后倾角
\(\alpha\)。将汽车等效为前、后各一轮，轴距为
\(d\)，质心位于轴距中点，前轮可转向、后轮不可转向，轮半径为
\(r\)。汽车质量为 \(m\)、速度为 \(v\)，前轮偏转小角
\(\theta\)，求侧向力造成的自动回正力矩。

\begin{figure}[H]
\centering
\includegraphics[width=0.9\linewidth,height=\textheight,keepaspectratio,alt={图1-1　主销内倾、主销后倾与轮胎接地点}]{../figures/question-01.pdf}
\caption*{图1-1　主销内倾、主销后倾与轮胎接地点}
\end{figure}

\newpage

\section{第 2 题}\label{ux7b2c-2-ux9898}

两个完全相同的永磁体相对放置，磁场的轴向分量为 \(B_z(z)\)；\(B_z\) 在
\(z=0\) 处最小，在 \(z=\pm L\) 处最大。带电粒子沿 \(z\)
方向的运动远慢于其在 \(xOy\) 平面内的回旋运动。粒子质量为 \(m\)、电荷为
\(q\)，从原点释放，初速度分量为 \(V_z(0)\) 与 \(V_\perp(0)\)。

（1）证明 \(V_\perp^2\) 与 \(B_z(z)\) 成正比；

（2）定义 \[
U(z)=\frac{mV_\perp^2(0)B_z(z)}{2B_z(0)},
\] 证明 \(\frac12mV_z^2+U(z)\) 为常量；

（3）在 \(z=\pm L\) 放置吸收板，定义
\(M=B_z(\pm L)/B_z(0)\)，求粒子不碰板的初速度条件；

（4）用初速度与 \(xOy\) 平面的夹角表示吸收条件。

\begin{figure}[H]
\centering
\includegraphics[width=0.847\linewidth,height=\textheight,keepaspectratio,alt={图2-1　磁镜场、回旋运动与两端吸收板}]{../figures/question-02.pdf}
\caption*{图2-1　磁镜场、回旋运动与两端吸收板}
\end{figure}

\newpage

\section{第 3 题}\label{ux7b2c-3-ux9898}

剪切模量定义为 \[
S=\frac{F/A}{\Delta L/L}=\frac{FL}{A\Delta L}.
\]

（1）密度为 \(\rho\)、剪切模量为 \(S\)、截面积为 \(A\)、长度为 \(L\)
的正方柱状``类果冻''物体，底面固定、顶面自由，仅沿正方形一边方向发生剪切振动。求其各阶振动周期。

（2）同样材料制成圆柱，截面积为 \(A\)、长度为
\(L\)，两端无外加力矩。沿截面圆周切向微扰后仅发生扭转振动，求各阶非零模态周期。

\begin{figure}[H]
\centering
\includegraphics[width=0.9\linewidth,height=\textheight,keepaspectratio,alt={图3-1　正方柱剪切振动与圆柱扭转振动}]{../figures/question-03.pdf}
\caption*{图3-1　正方柱剪切振动与圆柱扭转振动}
\end{figure}

\newpage

\section{第 4 题}\label{ux7b2c-4-ux9898}

在气液共存的等温线上，液态体积为 \(V_1\)、气态体积为 \(V_2\)，温度为
\(T\) 时的饱和蒸气压为 \(p(T)\)。在 \(T\) 与 \(T+\mathrm dT\)
两条等温线之间构造微小卡诺循环，证明 \[
\mathrm dp\,(V_2-V_1)=\mathrm dT\,(S_2-S_1).
\]

超导---正常相变中，正常态磁矩 \(M=0\)，超导态满足
\(M=-\mu_0B\)；外场对磁矩做功为 \(\mathrm dW=-B\,\mathrm dM\)。临界场 \[
B_c(T)=B_0\left(1-\frac{T^2}{T_0^2}\right).
\] 设 \(S_n(T)\)、\(S_s(T)\) 分别为正常态和超导态的熵，求
\(\Delta S_n(T)=S_n-S_s\)，并用 \[
B_0=0.144{\rm\,T},\quad
\mu_0=1.25\times10^{-6}{\rm\,N/A^2},\quad
T=1.95{\rm\,K},\quad T_0=4.48{\rm\,K}
\] 计算数值。

\begin{figure}[H]
\centering
\includegraphics[width=0.849\linewidth,height=\textheight,keepaspectratio,alt={图4-1　气液相变的微小卡诺循环}]{../figures/question-04.pdf}
\caption*{图4-1　气液相变的微小卡诺循环}
\end{figure}

\newpage

\section{第 5 题}\label{ux7b2c-5-ux9898}

在 \(xOy\) 平面内有沿 \(+z\) 方向的非均匀磁场 \[
B(x,y)=B_0\cos kx\cos ky.
\] 平面内有一个自感为 \(L\)、面积为 \(S\)、质量为 \(m\)
的很小刚性线圈，形状任意，初始位于原点，电流为
\(I_0>0\)，逆时针为正，电阻忽略。

（1）线圈移动到 \((x,y)\) 时，求电流 \(I(x,y)\)；

（2）求此时线圈在平面内的受力；

（3）线圈在原点受到 \(x\) 方向微扰，求小振动规律。

\begin{figure}[H]
\centering
\includegraphics[width=0.648\linewidth,height=\textheight,keepaspectratio,alt={图5-1　小线圈在二维周期磁场中的位置}]{../figures/question-05.pdf}
\caption*{图5-1　小线圈在二维周期磁场中的位置}
\end{figure}

\newpage

\section{第 6 题}\label{ux7b2c-6-ux9898}

由相同电阻构成一端封闭、向另一端无限延伸的方形管状电阻网络，每条边的电阻均为
\(r\)。端面四个顶点中，\(A\) 与 \(B\) 为对角点，\(A\) 与 \(C\)
为相邻点。求等效电阻 \(r_{AB}\) 与 \(r_{AC}\)。

\begin{figure}[H]
\centering
\includegraphics[width=0.9\linewidth,height=\textheight,keepaspectratio,alt={图6-1　半无限方形管状电阻网络}]{../figures/question-06.pdf}
\caption*{图6-1　半无限方形管状电阻网络}
\end{figure}

\newpage

\section{第 7 题}\label{ux7b2c-7-ux9898}

质量为 \(m\) 的相对论性粒子在一维偶势 \(V(x)\)
中运动。势能在原点取最小值0，并在 \(x>0\) 单调增大。粒子在原点的动能为
\(E_{k0}=V(b)\)，故振幅为 \(b\)。

（1）写出速度 \(v(x)\)；

（2）写出周期的定积分；

（3）取 \(V(x)=\frac12kx^2\)，定义
\(a=k/(2mc^2)\)，保留最低阶相对论修正，用 \(a,b\) 表示周期；

（4）求修正项与非相对论主项之比，并用 \(\eta=E_{k0}/(mc^2)\) 表示。

\begin{figure}[H]
\centering
\includegraphics[width=0.632\linewidth,height=\textheight,keepaspectratio,alt={图7-1　相对论性谐振子的能量与转向点}]{../figures/question-07.pdf}
\caption*{图7-1　相对论性谐振子的能量与转向点}
\end{figure}

\newpage

\section{第 8 题}\label{ux7b2c-8-ux9898}

有限光束可看作中心频率或中心入射角附近的平面波叠加。

（1）两列波 \[
U_1=U_0\cos(kx-\omega t),\qquad
U_2=U_0\cos[(k+\mathrm dk)x-(\omega+\mathrm d\omega)t]
\] 沿 \(x\) 轴传播，写出叠加波，并求相速度与群速度。

（2）频率为 \(\omega\) 的光从折射率 \(n_1\) 的介质以角度 \(\theta\)
入射到折射率 \(n_2\) 的介质。用 \(\theta\) 与 \(\theta+\mathrm d\theta\)
两束波模拟角谱，在界面 \(x\) 轴写出叠加波。

（3）全反射时 \[
\sin\theta_2=\frac{n_1\sin\theta}{n_2}>1,\qquad
\cos\theta_2=-i\sqrt{\left(\frac{n_1\sin\theta}{n_2}\right)^2-1},
\] 垂直偏振的振幅反射系数为 \[
r=\frac{n_1\cos\theta-n_2\cos\theta_2}
{n_1\cos\theta+n_2\cos\theta_2}.
\] 求反射相位差 \(\delta(\theta)\)。

（4）求反射波包中心相对几何反射点的横向位移 \(\Delta(\theta)\)。

\begin{figure}[H]
\centering
\includegraphics[width=0.88\linewidth,height=\textheight,keepaspectratio,alt={图8-1　角谱叠加、全反射与界面位移}]{../figures/question-08.pdf}
\caption*{图8-1　角谱叠加、全反射与界面位移}
\end{figure}
\end{document}
